\documentclass[11pt]{article}
\usepackage[T1]{fontenc}
\usepackage[utf8]{inputenc}
\usepackage{amsmath,amssymb,amsthm,mathtools}
\usepackage{booktabs}
\usepackage{geometry}
\usepackage{microtype}
\usepackage[hidelinks]{hyperref}
\geometry{margin=1in}

\newtheorem{theorem}{Theorem}
\newtheorem{proposition}[theorem]{Proposition}
\newtheorem{corollary}[theorem]{Corollary}
\newtheorem{lemma}[theorem]{Lemma}
\theoremstyle{definition}
\newtheorem{definition}[theorem]{Definition}
\newtheorem{remark}[theorem]{Remark}

\newcommand{\Torus}{\mathbb T^3}
\newcommand{\PN}{P_N}
\newcommand{\QN}{Q_N}
\newcommand{\eps}{\varepsilon}

\hypersetup{
  pdftitle={From Energy Observations to Continuum Error Certificates for Navier--Stokes},
  pdfauthor={Yaoharee Lahtee},
  pdfsubject={energy observability; discrete epsilon-completion; certified finite-to-continuum composition},
  pdfkeywords={Navier--Stokes, observability, Fourier-Galerkin, epsilon-completion, certification, relative energy}
}

\title{From Energy Observations to Continuum Error Certificates\\
for Finite Fourier--Galerkin Navier--Stokes}
\author{Yaoharee Lahtee\\
\small Open Civil Science Initiative, Bangkok, Thailand}
\date{September 2026}

\begin{document}
\maketitle

\begin{abstract}
Two different completeness questions arise in finite numerical fluid dynamics.  The first is
\emph{inner}: whether a declared observation process contains enough information to determine the
retained finite state, up to unavoidable symmetry.  The second is \emph{outer}: whether the
unrepresented continuum tail is provably small.  A companion finite observability analysis of
periodic three-dimensional Fourier--Galerkin Navier--Stokes gives all-resolution structural rank
ceilings for total-energy and shell-energy Lie jets, positive-viscosity rank universality, and exact
ceiling saturation in four certified reader--resolution cases.  Separately, the Discrete
Epsilon-Completion (EPSC) programme supplies target-indexed omitted-tail certificates, including
Leray--Hopf spacetime bounds and an a posteriori residual-based terminal adapter with an exact-dyadic
continuous-time reconstruction of recorded RK4 nodes.  This paper composes the two layers.  If an
observation/inversion layer supplies a certified retained-state reconstruction radius $\rho_N$
modulo an isometric symmetry preserving the Fourier cutoff, and EPSC supplies a certified omitted
tail radius $\beta_N$, orthogonality gives
\[
\inf_{g\in G}\|u(T)-g\widehat x_N\|_2
\le \sqrt{\rho_N^2+\beta_N^2}.
\]
The result makes precise the distinction between local observability rank and a quantitative
measurement-to-continuum certificate.  The missing inner-side ingredient is a certified inversion
radius from energy jets; stable recovery under noisy high-order time information is a further open
problem.  These frontiers are distinct from the EPSC high-cutoff tightness/scalability problem.
No claim about global regularity, finite-time singularity, global injectivity, or the Clay
Millennium problem is made.
\end{abstract}

\section{Introduction}

A numerical state may be incomplete for two logically independent reasons.  First, the sensors or
readouts may not determine the finite state that the solver actually retains.  Second, even a fully
known finite state may omit continuum information above the cutoff.  Treating these as one problem
obscures both the mathematics and the engineering design space.

We therefore write the certification chain as
\[
\boxed{
\text{measurements}
\longrightarrow
\text{retained-state certificate}
\longrightarrow
\text{continuum-tail certificate}
\longrightarrow
\text{final tolerance verdict}.
}
\]
The first arrow is an observability/inversion problem.  The second is an EPSC problem.  The present
paper states their exact composition rule and identifies which parts are already established and
which remain open.

The companion manuscript \emph{Energy Observability Across Fourier--Galerkin Resolutions of
Three-Dimensional Navier--Stokes} develops the finite observability side.  The EPSC manuscript
\emph{Finite Readout Certificates for Navier--Stokes} develops the finite-to-continuum side.  The
purpose here is synthesis, not duplication.

\section{Finite Fourier--Galerkin state and energy readers}

For a cubic cutoff
\[
K_N=\{-N,\ldots,N\}^3\setminus\{0\},
\]
the declared real phase-space dimension is
\begin{equation}\label{eq:dN}
\boxed{d_N=2\big((2N+1)^3-1\big).}
\end{equation}
Let $Q_N^{\rm shell}=\{|k|^2:k\in K_N\}$ and $m_N=|Q_N^{\rm shell}|$.  Let $I$ be the vector of
shell energies and $E=\mathbf 1^T I$ the total retained kinetic energy.  For the finite vector field
$F_\nu$, define the scalar and shell Lie jets
\[
\mathcal E_R=(E,\mathcal L_FE,\ldots,\mathcal L_F^R E),
\qquad
\mathcal J_R=(I,\mathcal L_FI,\ldots,\mathcal L_F^R I).
\]
Spatial translation acts on Fourier coefficients by
\[
(T_a\widehat u)_k=e^{ik\cdot a}\widehat u_k.
\]
The dynamics is translation-equivariant while both energy readers are invariant.  At generic states
where the translation action is locally free, no invariant observation map can exceed rank $d_N-3$.

The finite observability analysis gives the structural ceilings
\begin{equation}\label{eq:scalarceiling}
\boxed{\operatorname{rank}D\mathcal E_R\le\min(R+1,d_N-3),}
\end{equation}
\begin{equation}\label{eq:shellceiling}
\boxed{\operatorname{rank}D\mathcal J_R\le
\min\big(m_N+(m_N-1)R,d_N-3\big).}
\end{equation}
Hence the symmetry ceiling cannot be attained before
\begin{equation}\label{eq:depths}
R_E^{\min}=d_N-4,
\qquad
R_I^{\min}=\left\lceil\frac{d_N-3-m_N}{m_N-1}\right\rceil.
\end{equation}

For every positive viscosity, the scaling $x=\nu y$, $s=\nu t$ gives
\[
\mathcal L_{F_\nu}^r h(\nu y)=\nu^{r+2}\mathcal L_{F_1}^r h(y),
\qquad
D_x\mathcal L_{F_\nu}^r h(\nu y)=\nu^{r+1}D_y\mathcal L_{F_1}^r h(y),
\]
for the quadratic energy readers.  The resulting jet ranks are therefore independent of the value
of $\nu>0$.  The Euler endpoint $\nu=0$ is separate.

\section{What is actually certified on the observability side}

Exact modular certificates currently attain the translation ceiling at the earliest structurally
admissible order in the following four cases:
\begin{center}
\begin{tabular}{ccrrr}
\toprule
$N$ & reader & $m_N$ & $R$ & certified rank \\
\midrule
1 & total energy & 1 & 48 & 49 \\
1 & shell energies & 3 & 23 & 49 \\
2 & total energy & 1 & 244 & 245 \\
3 & shell energies & 18 & 39 & 681 \\
\bottomrule
\end{tabular}
\end{center}
Every row reaches $d_N-3$.  At a generic free-action state, this means that the kernel of the
saturated jet differential is exactly the tangent space of the translation orbit.  It is therefore
appropriate to call the result \emph{local finite-state completeness modulo translations}.

However, three stronger statements do not follow from rank saturation: global injectivity of the
energy jet; a constructive inverse with branch control; or a numerical reconstruction radius under
measurement error.  The all-resolution assertion that the ceilings are generically attained at the
earliest admissible order remains a conjecture.

The structural formulas also expose a measurement--time tradeoff.  One scalar channel requires at
least $d_N-4=\Theta(N^3)$ time derivatives before saturation is even possible.  A multi-shell reader
can reduce this depth.  At $N=3$, the scalar lower bound is $R\ge680$, whereas the 18-channel shell
reader is exactly saturated at $R=39$.  This is a structural observability statement, not yet a
noise-robust sensor theorem.

\section{The EPSC outer-completeness layer}

Let $P_N$ be the retained Fourier projection and $Q_N=I-P_N$.  EPSC distinguishes nested finite
agreement from a proved omitted-tail bound.  A target-indexed certificate has the form
\[
\|Q_Nu\|_Y\le\beta_{N,Y}.
\]
For an unforced Leray--Hopf solution on the $2\pi$-periodic three-torus,
\begin{equation}\label{eq:stbeta}
\boxed{
\|Q_Nu\|_{L^2(0,T;L^2_x)}
\le
\frac{\|u_0\|_2}{\sqrt{2\nu}(N+1)}.
}
\end{equation}
This is a genuine $O(N^{-1})$ spacetime tail certificate and does not assume global smoothness.

For a prescribed terminal time, richer retained information can yield a terminal bound.  One route
is the energy/dissipation budget; another uses a comparison path $v$ and the residual
\[
r=\partial_tv+\mathbb P[(v\cdot\nabla)v]-\nu\Delta v-\mathbb Pf.
\]
With
\[
A_T=2\int_0^T\|\nabla v\|_\infty dt,
\qquad
B_T=\int_0^T\|r\|_{H^{-1}}^2dt,
\]
the standard relative-energy estimate gives
\begin{equation}\label{eq:RE}
\sup_{t\le T}\|u(t)-v(t)\|_2^2
\le e^{A_T}\left(e_0^2+\frac{B_T}{\nu}\right).
\end{equation}
If $v(T)$ is $P_N$-supported, the same bound controls the terminal omitted tail.

The current EPSC implementation supplies a concrete continuous comparison path from recorded RK4
nodes: every binary64 component is interpreted as its exact dyadic rational, the nodes are projected
exactly onto divergence-free Fourier space, and consecutive nodes are joined by a rational
piecewise-linear path.  The resulting Navier--Stokes residual is polynomial on each time cell, so
its $H^{-1}$ squared integral can be evaluated exactly in rational arithmetic.  This closes the
finite-tape enclosure construction registered as \texttt{PROP-EPSC-15}.  Tightness and
computational scaling at large cutoff/horizon remain the distinct open \texttt{PROP-EPSC-16}.

\section{Observable-to-continuum composition}

The two layers can now be connected without identifying local rank with continuum accuracy.

\begin{theorem}[Observable-to-continuum orthogonal composition]\label{thm:compose}
Let $G$ be an isometric symmetry group on $L^2$ that preserves $P_N$.  Suppose a certified finite
observation/inversion procedure returns a $P_N$-supported state $\widehat x_N$ and radius $\rho_N$
such that
\begin{equation}\label{eq:rho}
\inf_{g\in G}\|P_Nu(T)-g\widehat x_N\|_2\le\rho_N.
\end{equation}
Suppose independently that an EPSC argument supplies
\begin{equation}\label{eq:betaN}
\|Q_Nu(T)\|_2\le\beta_N.
\end{equation}
Then
\begin{equation}\label{eq:compose}
\boxed{
\inf_{g\in G}\|u(T)-g\widehat x_N\|_2
\le
\sqrt{\rho_N^2+\beta_N^2}.
}
\end{equation}
\end{theorem}

\begin{proof}
For any $g\in G$ admissible in~\eqref{eq:rho}, the difference
$P_Nu(T)-g\widehat x_N$ lies in the retained subspace, while $Q_Nu(T)$ lies in its orthogonal
complement.  Because $u=P_Nu+Q_Nu$,
\[
\|u(T)-g\widehat x_N\|_2^2
=
\|P_Nu(T)-g\widehat x_N\|_2^2+\|Q_Nu(T)\|_2^2.
\]
Apply the two certified bounds and take the infimum over $g$.
\end{proof}

For the energy-observability application, $G=\Torus$ acts by spatial translations.  Thus a final
fail-closed tolerance rule is
\begin{equation}\label{eq:finalgate}
\boxed{
\sqrt{\rho_N^2+\beta_N^2}\le\eps
\quad\Longrightarrow\quad
\text{continuum state certified modulo translation at tolerance }\eps.
}
\end{equation}
If either $\rho_N$ or $\beta_N$ is absent, the verdict is \texttt{HOLD}.

\begin{remark}[Why this is sharper than addition]
A generic triangle inequality would give $\rho_N+\beta_N$.  The Fourier split is orthogonal, so the
retained and omitted errors combine Pythagoreanly.  This is useful because the observation-side and
tail-side budgets can be allocated independently under a single final tolerance.
\end{remark}

\section{Symmetry compatibility}

Total and shell energies are invariant under spatial translations, so finite observability can at
best recover a translation orbit.  This ambiguity is not a defect of the reader; it is the symmetry
of the model.  The continuum $L^2$ tail norm is also translation invariant, and $P_N$ commutes with
translations.  Therefore Theorem~\ref{thm:compose} naturally lives on the quotient by translations.

For residual-based EPSC certificates, the exact mathematical quantities $A_T$ and $B_T$ are
translation invariant when the comparison path and solution are translated together.  A specific
computer majorant may nevertheless depend on the chosen Fourier phase representation and can be
looser for one representative than another.  Certificate validity is preserved; tightness and
representative optimisation belong to the engineering/scalability frontier rather than to the
observability theorem.

\section{The three frontiers after composition}

The combined architecture separates three open tasks that should not be collapsed into one.

\paragraph{F1: all-resolution observability saturation.}
The structural ceilings~\eqref{eq:scalarceiling}--\eqref{eq:shellceiling} hold at every finite
resolution, but exact generic saturation at the earliest admissible depth is certified only in the
four recorded cases.  The all-$N$ statement remains \texttt{PROP-NSOBS-07}.

\paragraph{F2: certified finite inversion and noise stability.}
Rank saturation yields local differential identifiability modulo translations, but the composition
theorem needs a certified radius $\rho_N$.  Constructing such an inverse is
\texttt{PROP-EPSC-18}; propagating measurement and differentiation uncertainty through it is
\texttt{PROP-EPSC-19}.

\paragraph{F3: scalable outer certification.}
The exact-dyadic EPSC-15 path construction is a correctness route, but its conservative gradient
majorants and residual computation may become expensive or loose as $N$ and $T$ grow.  Improving
that cost/tightness is \texttt{PROP-EPSC-16}.

These tasks interact in an eventual practical system, but they are mathematically different.

\section{Executable evidence and provenance}

The repository reproduction ledger contains finite checks for the observability bookkeeping,
recorded saturation ceilings, translation invariance of the energy readers, the orthogonal
$\sqrt{\rho^2+\beta^2}$ composition, EPSC finite algebra, and the exact-dyadic continuous-time RK4
path enclosure.  Standard analytic implications remain labelled at the analytic/derived tier rather
than being presented as machine proofs.

Proposal provenance is split deliberately.  The NS observability family is registered as
\texttt{PROP-NSOBS-01} through \texttt{PROP-NSOBS-08}.  The EPSC composition extension is registered
as \texttt{PROP-EPSC-17} through \texttt{PROP-EPSC-19}; \texttt{PROP-EPSC-16} remains the separate
high-cutoff tightness/scalability frontier.  Proposal identifiers are not automatically canonical
Toledo theorem codes.

\section{Claim boundary}

The strongest new statement here is a conditional composition theorem: a certified finite
reconstruction radius and a certified omitted-tail radius combine orthogonally into a continuum
quotient-state radius.  This does not turn qualitative local rank into a numerical inverse, and it
does not prove that total or shell energy observations are globally injective.  The exact finite
saturation results remain finite-resolution observability results.  The EPSC certificates remain
target- and assumption-specific.

Nothing in this paper proves global smoothness, finite-time blow-up, uniqueness of all Leray--Hopf
weak solutions, continuum observability from finitely many noisy sensors, turbulent DNS adequacy at
coarse cutoff, or the Clay Millennium existence-and-smoothness problem.

\section{Conclusion}

Energy observability and EPSC close opposite sides of the same measurement-to-continuum pipeline.
The former asks whether the represented state is visible; the latter asks how much unrepresented
state can remain.  Their rigorous interface is not ``rank plus convergence looks good'', but two
certified radii.  Once a finite inversion radius $\rho_N$ and an omitted-tail radius $\beta_N$ are
available, the final $L^2$ quotient-state certificate is
\[
\sqrt{\rho_N^2+\beta_N^2}.
\]
This cleanly separates what has been established from what remains: all-resolution saturation,
quantitative/noise-stable energy-jet inversion, and scalable tight outer certification.

\begin{thebibliography}{9}
\bibitem{Leray1934}
J. Leray,
\newblock Sur le mouvement d'un liquide visqueux emplissant l'espace,
\newblock \emph{Acta Mathematica} 63 (1934), 193--248.

\bibitem{Hopf1951}
E. Hopf,
\newblock \"Uber die Anfangswertaufgabe f\"ur die hydrodynamischen Grundgleichungen,
\newblock \emph{Mathematische Nachrichten} 4 (1951), 213--231.

\bibitem{Temam}
R. Temam,
\newblock \emph{Navier--Stokes Equations: Theory and Numerical Analysis},
\newblock AMS Chelsea Publishing.

\bibitem{Fefferman}
C. L. Fefferman,
\newblock Existence and smoothness of the Navier--Stokes equation,
\newblock Clay Mathematics Institute Millennium Problem description, 2000.

\bibitem{LahteeReadout}
Y. Lahtee,
\newblock A Readout Problem for Fefferman's Existence and Smoothness of the Navier--Stokes Equation,
\newblock Zenodo, DOI: 10.5281/zenodo.22673246, 2026.

\bibitem{LahteeObs}
Y. Lahtee,
\newblock Energy Observability Across Fourier--Galerkin Resolutions of Three-Dimensional Navier--Stokes,
\newblock companion manuscript and executable record, 2026.

\bibitem{LahteeEPSC}
Y. Lahtee,
\newblock Finite Readout Certificates for Navier--Stokes: Discrete Epsilon-Completion, Spectral Tails,
\newblock and A Posteriori Adapters, companion manuscript and executable record, 2026.
\end{thebibliography}

\end{document}
